% ============================================================================
% IEEE Conference Paper: M.A.T.R.I.X. AI Operating System
% ============================================================================
% Compile with: pdflatex -> bibtex -> pdflatex -> pdflatex
% ============================================================================

\documentclass[conference]{IEEEtran}

% ── Packages ────────────────────────────────────────────────────────────────
\usepackage{cite}
\usepackage{amsmath,amssymb,amsfonts}
\usepackage{algorithmic}
\usepackage{algorithm}
\usepackage{graphicx}
\usepackage{textcomp}
\usepackage{xcolor}
\usepackage{hyperref}
\usepackage{booktabs}
\usepackage{multirow}
\usepackage{array}
\usepackage{enumitem}
\usepackage{balance}
\usepackage{url}
\usepackage{listings}
\usepackage{subcaption}

% ── Listings Configuration ──────────────────────────────────────────────────
\lstset{
  basicstyle=\ttfamily\scriptsize,
  breaklines=true,
  frame=single,
  numbers=left,
  numberstyle=\tiny\color{gray},
  keywordstyle=\color{blue}\bfseries,
  commentstyle=\color{green!50!black},
  stringstyle=\color{red!70!black},
  tabsize=2,
  showstringspaces=false
}

% ── Metadata ────────────────────────────────────────────────────────────────
\def\BibTeX{{\rm B\kern-.05em{\sc i\kern-.025em b}\kern-.08em
    T\kern-.1667em\lower.7ex\hbox{E}\kern-.125emX}}

\begin{document}

% ============================================================================
%                              TITLE & AUTHORS
% ============================================================================

\title{M.A.T.R.I.X.: A Hardware-Adaptive, Offline-First AI Operating System with Multi-Policy Agent Scheduling and Defense-in-Depth Security}

\author{
\IEEEauthorblockN{Tauhid Ahmed}
\IEEEauthorblockA{
Department of Electrical and Electronic Engineering \\
Independent Researcher \\
\textit{tauhid@matrix-os.dev}
}
}

\maketitle

% ============================================================================
%                               ABSTRACT
% ============================================================================

\begin{abstract}
The proliferation of large language models (LLMs) has catalyzed a paradigm shift in human-computer interaction, yet contemporary AI assistants remain tethered to cloud infrastructure, raising critical concerns regarding privacy, latency, and availability. This paper presents \textbf{M.A.T.R.I.X.} (\textit{Metaverse Artificial Technological Regenerating Intelligence eXperiment}), a novel custom Linux-based operating system that integrates local AI inference, multi-policy agent task scheduling, hardware-adaptive model selection, and defense-in-depth security into a unified, offline-first computing environment. The system implements four classical CPU scheduling algorithms---FIFO, Shortest Job First (SJF), Round Robin, and Priority with Aging---augmented by an LRU-K attention cache for efficient agent context switching. A hardware profiling engine automatically classifies target machines into four performance tiers and selects optimal LLM configurations ranging from 1.5 billion to 70 billion parameters. Security is enforced through a four-layer architecture comprising AppArmor Mandatory Access Control, Bubblewrap namespace isolation, systemd service hardening, and application-level command classification. The operating system is compiled from a 13-stage automated ISO build pipeline based on Debian Bookworm, validated through headless QEMU boot verification, and managed via a React-based glassmorphic desktop interface featuring voice-controlled AI interaction, real-time hardware telemetry, and a 3D holographic visualization core. The entire backend is implemented using exclusively Python standard library modules with zero external dependencies, achieving a total codebase of 3,283 lines across 17 source files. Experimental evaluation demonstrates sub-200ms context switch latency, effective starvation prevention through priority aging, and successful automated boot-to-service initialization within 60 seconds on commodity hardware. M.A.T.R.I.X. represents a comprehensive proof-of-concept for privacy-preserving, self-configuring AI operating systems capable of functioning entirely without network connectivity.
\end{abstract}

\begin{IEEEkeywords}
Artificial Intelligence Operating System, Local LLM Inference, Agent Task Scheduling, LRU-K Cache, Hardware-Adaptive Computing, Defense-in-Depth Security, Offline-First Architecture, AppArmor, Bubblewrap Sandboxing
\end{IEEEkeywords}

% ============================================================================
%                          I. INTRODUCTION
% ============================================================================

\section{Introduction}

The integration of artificial intelligence into operating systems represents a fundamental evolution in computing paradigms. While modern AI assistants such as ChatGPT \cite{openai2023gpt4}, Google Gemini \cite{google2024gemini}, and Microsoft Copilot \cite{microsoft2024copilot} have demonstrated remarkable capabilities in natural language understanding, code generation, and task automation, they remain fundamentally dependent on cloud infrastructure. This dependency introduces three critical limitations: (1)~\textit{privacy vulnerability}, where sensitive user data traverses external servers; (2)~\textit{latency constraints}, where network round-trip times degrade real-time interaction quality; and (3)~\textit{availability fragility}, where network outages render the AI assistant entirely non-functional.

Concurrently, the open-source LLM ecosystem has matured significantly, with models such as Llama~3 \cite{meta2024llama3}, Qwen~2.5 \cite{qwen2024qwen25}, and Phi-3 \cite{microsoft2024phi3} achieving competitive performance at parameter counts feasible for local inference on consumer hardware. Runtime frameworks including Ollama \cite{ollama2024} and llama.cpp \cite{llamacpp2024} have further reduced the computational barrier, enabling real-time text generation on CPUs with as little as 4\,GB of system RAM.

Despite these advances, no existing operating system provides a holistic integration of local AI inference with operating system-level resource management, security sandboxing, and hardware-adaptive optimization. Existing approaches either bolt AI capabilities onto traditional OS shells (e.g., Windows Copilot) or provide standalone AI applications without OS-level integration (e.g., LM~Studio, GPT4All).

This paper presents \textbf{M.A.T.R.I.X.} (\textit{Metaverse Artificial Technological Regenerating Intelligence eXperiment}), a custom Linux operating system designed from the ground up to address these gaps. The principal contributions of this work are:

\begin{enumerate}[label=(\arabic*)]
    \item A \textbf{hardware-adaptive AI model selection engine} that automatically profiles system capabilities and selects optimal LLM configurations across four performance tiers, from 1.5B to 70B parameter models.
    
    \item A \textbf{multi-policy agent task scheduler} implementing four classical scheduling algorithms (FIFO, SJF, Round Robin, Priority with Aging) augmented by an LRU-K attention cache for minimizing agent context switch overhead.
    
    \item A \textbf{four-layer defense-in-depth security architecture} combining AppArmor MAC profiles, Bubblewrap namespace containers, systemd service hardening, and application-level permission classification.
    
    \item A \textbf{zero-dependency backend architecture} where all Python components use exclusively standard library modules, eliminating external package management within the minimal OS environment.
    
    \item A \textbf{premium web-based desktop interface} built with React~19 featuring glassmorphic design, voice-controlled AI interaction via Web Speech API, 3D holographic visualization, and real-time hardware telemetry.
    
    \item A \textbf{fully automated ISO compilation pipeline} with 13 build stages, headless QEMU boot verification, and CI/CD integration through a multi-agent GitHub Actions workflow chain.
\end{enumerate}

The remainder of this paper is organized as follows: Section~\ref{sec:literature} reviews related work across AI-integrated operating systems, scheduling algorithms, and OS security mechanisms. Section~\ref{sec:architecture} presents the overall system architecture. Sections~\ref{sec:scheduler}--\ref{sec:security} detail the scheduling subsystem, hardware profiler, and security model. Section~\ref{sec:ui} describes the user interface. Section~\ref{sec:build} covers the ISO build pipeline. Section~\ref{sec:results} presents experimental evaluation. Section~\ref{sec:future} discusses future work, and Section~\ref{sec:conclusion} concludes the paper.

% ============================================================================
%                      II. LITERATURE REVIEW
% ============================================================================

\section{Literature Review}
\label{sec:literature}

\subsection{AI-Integrated Operating Systems}

The concept of embedding intelligence into operating systems has a long history. Early expert system shells such as MYCIN \cite{shortliffe1976mycin} and CLIPS \cite{giarratano2005expert} demonstrated rule-based reasoning within domain-specific applications but lacked OS-level integration. Microsoft's Cortana \cite{microsoft2020cortana} and Apple's Siri \cite{apple2024siri} introduced voice-activated AI assistants at the OS level, but these systems fundamentally rely on cloud processing and do not perform local inference.

The emergence of on-device LLMs has opened new possibilities. Google's Gemini Nano \cite{google2024gemini} runs inference on Pixel devices, and Apple Intelligence \cite{apple2024intelligence} performs local processing on Apple Silicon. However, these implementations are proprietary, closed-source, and tightly coupled to specific hardware ecosystems.

In the open-source domain, projects such as GPT4All \cite{gpt4all2024} and LM~Studio \cite{lmstudio2024} provide desktop applications for local LLM inference, but they operate as standalone applications rather than OS-integrated services. Jan.ai \cite{jan2024} offers a self-hosted AI platform but lacks operating system-level scheduling and resource management. NixOS-based AI configurations \cite{nixos2024} demonstrate declarative OS builds with AI tooling but do not implement custom scheduling or security sandboxing.

M.A.T.R.I.X. distinguishes itself by providing complete OS-level integration where the AI inference engine, task scheduler, security layer, and user interface are architecturally unified within a custom Linux distribution, rather than functioning as separate applications running atop a general-purpose OS.

\subsection{CPU Scheduling Algorithms}

Process scheduling remains a foundational topic in operating system design \cite{silberschatz2018os, tanenbaum2015modern}. The four algorithms implemented in M.A.T.R.I.X.---FIFO, SJF, Round Robin, and Priority---are well-established in classical OS literature.

\textbf{FIFO} (First-In-First-Out) provides the simplest scheduling policy with $O(1)$ selection complexity using queue data structures \cite{stallings2018os}. \textbf{SJF} (Shortest Job First) is provably optimal for minimizing average waiting time in non-preemptive scenarios \cite{silberschatz2018os}, though it requires burst time estimation. \textbf{Round Robin} introduces time-slice preemption, offering fairness guarantees at the cost of increased context switch overhead \cite{tanenbaum2015modern}. \textbf{Priority scheduling with aging} addresses the starvation problem inherent in static priority schemes by dynamically incrementing the priority of waiting processes \cite{stallings2018os}.

The LRU-K page replacement algorithm, introduced by O'Neil et al. \cite{oneil1993lruk}, extends traditional LRU by considering the $K$-th backward reference distance rather than only the most recent access. This approach has been widely adopted in database buffer management but has seen limited application in agent context caching. M.A.T.R.I.X. adapts LRU-K to track AI agent ``attention history'' states, representing a novel application domain for this algorithm.

\subsection{Operating System Security Mechanisms}

Modern Linux security employs a defense-in-depth strategy combining multiple containment layers. AppArmor \cite{apparmor2024} provides path-based Mandatory Access Control (MAC), restricting program capabilities to defined file and network access patterns. Bubblewrap \cite{bubblewrap2024} implements lightweight sandboxing through Linux namespace isolation (PID, network, mount, IPC, UTS), providing container-like confinement without the overhead of full virtualization. Systemd service hardening \cite{systemd2024} offers process-level sandboxing through capabilities bounding, filesystem protection, and privilege restriction directives.

Flatpak \cite{flatpak2024}, Snap \cite{snap2024}, and Firejail \cite{firejail2024} provide application-level sandboxing for desktop Linux, but these are package-level solutions rather than OS-level security architectures. M.A.T.R.I.X. integrates all four layers---AppArmor, Bubblewrap, systemd, and application-level classification---into a coherent security model specifically designed for AI agent containment.

\subsection{Custom Linux Distribution Engineering}

The construction of custom Linux distributions from minimal bases has been extensively documented. Debian Live \cite{debianlive2024} provides the \texttt{live-build} framework, while the Linux From Scratch project \cite{lfs2024} offers manual compilation guidance. Tools such as \texttt{debootstrap} \cite{debootstrap2024} enable minimal chroot bootstrapping, and SquashFS \cite{squashfs2024} provides compressed read-only filesystem images for live environments.

Automated OS deployment through preseed files \cite{debianpreseed2024} and unattended installation frameworks reduces manual intervention. QEMU-based validation \cite{qemu2024} enables headless boot testing for CI/CD integration. M.A.T.R.I.X. synthesizes these established tools into an integrated 13-stage compilation pipeline with automated verification.

% ============================================================================
%                    III. SYSTEM ARCHITECTURE
% ============================================================================

\section{System Architecture}
\label{sec:architecture}

M.A.T.R.I.X. employs a layered architecture comprising four principal subsystems, as illustrated in Fig.~\ref{fig:architecture}:

\begin{figure}[htbp]
\centering
\setlength{\unitlength}{1cm}
\begin{picture}(8.5,9.5)
% Layer 4 - UI
\put(0.25,8.0){\framebox(8.0,1.2){}}
\put(0.5,8.7){\small\textbf{Layer 4: Operator Interface (React 19 + Vite 6)}}
\put(0.5,8.2){\scriptsize Terminal | Build Factory | Telemetry | Sandbox | AI Voice | File Explorer}

% Layer 3 - API
\put(0.25,6.2){\framebox(8.0,1.4){}}
\put(0.5,7.2){\small\textbf{Layer 3: Unified API Daemon (Python HTTP)}}
\put(0.5,6.7){\scriptsize REST Endpoints | Token Auth | CORS | Command Classification}
\put(0.5,6.4){\scriptsize Audit Logging | Path Confinement | SPA Serving}

% Layer 2 - Core Services
\put(0.25,4.0){\framebox(3.8,1.8){}}
\put(0.5,5.4){\small\textbf{Scheduler Daemon}}
\put(0.5,5.0){\scriptsize FIFO | SJF | RR | Priority}
\put(0.5,4.6){\scriptsize LRU-K Cache}
\put(0.5,4.2){\scriptsize Ollama Integration}

\put(4.45,4.0){\framebox(3.8,1.8){}}
\put(4.7,5.4){\small\textbf{Hardware Profiler}}
\put(4.7,5.0){\scriptsize CPU/RAM/GPU Detection}
\put(4.7,4.6){\scriptsize Tier Classification}
\put(4.7,4.2){\scriptsize Model Recommendation}

% Layer 1 - Security
\put(0.25,2.0){\framebox(8.0,1.6){}}
\put(0.5,3.2){\small\textbf{Layer 1: Security Confinement}}
\put(0.5,2.7){\scriptsize AppArmor MAC | Bubblewrap Namespaces | Systemd Hardening}
\put(0.5,2.3){\scriptsize nftables Firewall | Command Classification | Audit Trail}

% Layer 0 - OS
\put(0.25,0.5){\framebox(8.0,1.1){}}
\put(0.5,1.2){\small\textbf{Layer 0: Debian Bookworm (Custom ISO)}}
\put(0.5,0.7){\scriptsize Linux Kernel 6.x | MATE Desktop | Ollama | Wine/DXVK}

% Arrows
\put(4.25,8.0){\vector(0,-1){0.4}}
\put(4.25,6.2){\vector(0,1){0.4}}
\put(4.25,6.2){\vector(0,-1){0.4}}
\put(2.25,4.0){\vector(0,-1){0.4}}
\put(6.45,4.0){\vector(0,-1){0.4}}
\put(4.25,2.0){\vector(0,-1){0.4}}

\end{picture}
\caption{M.A.T.R.I.X. layered system architecture showing the five principal layers from hardware (Layer~0) through the operator interface (Layer~4).}
\label{fig:architecture}
\end{figure}

\subsection{Design Principles}

The architecture adheres to five core design principles:

\begin{enumerate}[label=\textbf{P\arabic*}:]
    \item \textbf{Offline-First}: The entire system operates without network connectivity. AI inference, scheduling, and user interaction function identically in air-gapped environments.
    
    \item \textbf{Zero External Dependencies}: All backend Python components use exclusively standard library modules (\texttt{os}, \texttt{sys}, \texttt{json}, \texttt{sqlite3}, \texttt{asyncio}, \texttt{http.server}, \texttt{subprocess}), eliminating \texttt{pip} dependency management.
    
    \item \textbf{Hardware Adaptivity}: The system dynamically reconfigures AI model selection, memory optimization, and thread allocation based on detected hardware capabilities.
    
    \item \textbf{Defense-in-Depth}: Security enforcement occurs at four independent layers, ensuring that compromise of any single layer does not grant unrestricted system access.
    
    \item \textbf{Privilege Separation}: System services execute under the unprivileged \texttt{matrix} user account with minimal capabilities.
\end{enumerate}

\subsection{Component Interaction Model}

The system employs a client-server architecture where the React frontend communicates with the Python API daemon via RESTful HTTP endpoints. The API daemon serves dual roles: (1)~providing REST API endpoints under the \texttt{/api/} path prefix, and (2)~serving the compiled React SPA with \texttt{index.html} fallback routing for all non-API routes.

Authentication is implemented through a per-session UUID4 token generated at daemon startup and injected into the frontend via a dynamically served \texttt{/api/token.js} endpoint. All subsequent API requests require this token in the \texttt{X-Matrix-Token} HTTP header, preventing unauthorized access from external browser contexts.

% ============================================================================
%                IV. MULTI-POLICY AGENT TASK SCHEDULER
% ============================================================================

\section{Multi-Policy Agent Task Scheduler}
\label{sec:scheduler}

The scheduling subsystem manages AI agent task execution through four interchangeable algorithms, unified by an LRU-K attention cache that minimizes context switch overhead.

\subsection{Task Model}

Each agent task $T_i$ is defined by the tuple:
\begin{equation}
T_i = \langle a_i, n_i, d_i, r_i, C_i, U_i, R_i, P_i \rangle
\end{equation}
where $a_i$ is the agent identifier, $n_i$ is the task name, $d_i$ is the estimated duration, $r_i$ is the remaining execution time, and $\{C_i, U_i, R_i, P_i\}$ are the multi-factor priority components: complexity, urgency, resource requirement, and user-assigned priority, respectively.

The composite priority score is computed as:
\begin{equation}
\label{eq:priority}
\mathcal{P}(T_i) = 0.2 \cdot C_i + 0.4 \cdot U_i + 0.1 \cdot R_i + 0.3 \cdot P_i
\end{equation}
where each factor is normalized to $[0, 10]$.

\subsection{Scheduling Algorithms}

\subsubsection{FIFO Scheduler}
Tasks are executed in arrival order using a FIFO queue. Selection complexity is $O(1)$ with \texttt{collections.deque}. Each task runs to completion (non-preemptive), executing for its full remaining time $r_i$.

\subsubsection{Shortest Job First (SJF)}
Tasks are sorted by remaining execution time $r_i$ in ascending order. The scheduler selects $T^* = \arg\min_{T_i \in Q} r_i$. While optimal for average waiting time minimization, SJF introduces potential starvation of long-duration tasks.

\subsubsection{Round Robin}
A fixed time quantum $q = 1.0$\,s is allocated per scheduling cycle. If $r_i > q$ after execution, the task is re-enqueued:
\begin{equation}
r_i' = r_i - q, \quad T_i \rightarrow \text{enqueue}(Q)
\end{equation}

\subsubsection{Priority with Aging}
Tasks are dispatched in descending priority order per Eq.~\eqref{eq:priority}. To prevent starvation, an aging mechanism increases the urgency of all waiting tasks after each scheduling round:
\begin{equation}
U_j' = \min(U_j + \Delta_{\text{age}}, U_{\max}), \quad \forall T_j \in Q \setminus \{T^*\}
\end{equation}
where $\Delta_{\text{age}} = 0.8$ and $U_{\max} = 10.0$. This ensures that even low-priority tasks eventually accumulate sufficient urgency for dispatch.

\subsection{LRU-K Attention Cache}

Agent context switches require saving and restoring attention state---the accumulated history of agent interactions. Naive approaches load/save states from the SQLite database on every switch, incurring significant I/O latency.

M.A.T.R.I.X. employs the LRU-K algorithm \cite{oneil1993lruk} to cache agent attention states in memory. For each cached agent $a$, the cache maintains an access history $H_a = [t_1, t_2, \ldots, t_n]$ recording timestamps of all accesses. The backward $K$-distance is defined as:
\begin{equation}
B_K(a) = t_{\text{current}} - H_a[|H_a| - K]
\end{equation}

Eviction selects the agent with the maximum $B_K$:
\begin{equation}
a^* = \arg\max_{a \in \text{Cache}} B_K(a)
\end{equation}

Agents with fewer than $K$ accesses are assigned $B_K = \infty$ and evicted first, ordered by earliest first access. In M.A.T.R.I.X., $K=2$ and cache capacity $=2$, yielding an LRU-2 policy.

\begin{algorithm}[htbp]
\caption{LRU-K Context Switch}
\label{alg:lruk}
\begin{algorithmic}[1]
\REQUIRE Agent $a_{\text{new}}$, Cache $\mathcal{C}$, Database $\mathcal{D}$
\STATE Save current agent state to $\mathcal{C}$ \COMMENT{50ms}
\IF{$a_{\text{new}} \in \mathcal{C}$}
    \STATE Load state from $\mathcal{C}$ \COMMENT{Cache HIT, 50ms}
    \STATE Update $H_{a_{\text{new}}}$ with current timestamp
\ELSE
    \STATE Load state from $\mathcal{D}$ \COMMENT{Cache MISS, 150ms}
    \IF{$|\mathcal{C}| \geq \text{capacity}$}
        \STATE $a^* \leftarrow \arg\max_{a \in \mathcal{C}} B_K(a)$
        \STATE Evict $a^*$ from $\mathcal{C}$
    \ENDIF
    \STATE Insert $a_{\text{new}}$ into $\mathcal{C}$
\ENDIF
\end{algorithmic}
\end{algorithm}

The cache hit path incurs 100\,ms total latency (50\,ms save + 50\,ms load) versus 200\,ms for a cache miss (50\,ms save + 150\,ms database load), representing a 50\% latency reduction on cache hits.

\subsection{Ollama LLM Integration}

The scheduler daemon optionally queries a local Ollama instance (\texttt{localhost:11434}) for dynamic task progress narration. Auto-discovery polls the \texttt{/api/tags} endpoint, prioritizing Qwen3 models. Requests use \texttt{asyncio.to\_thread()} for non-blocking HTTP communication with 2-second connection timeouts and 8-second generation timeouts.

% ============================================================================
%              V. HARDWARE-ADAPTIVE AI MODEL SELECTION
% ============================================================================

\section{Hardware-Adaptive AI Model Selection}
\label{sec:hardware}

The hardware profiler engine executes during first boot to classify the target system and configure optimal AI parameters.

\subsection{Detection Pipeline}

The profiler queries three hardware dimensions through platform-specific methods:

\begin{itemize}
    \item \textbf{CPU Cores}: \texttt{os.cpu\_count()} (cross-platform)
    \item \textbf{System RAM}: \texttt{/proc/meminfo} (Linux), \texttt{GlobalMemoryStatusEx} via \texttt{ctypes} (Windows), \texttt{sysctl hw.memsize} (macOS)
    \item \textbf{GPU VRAM}: \texttt{nvidia-smi --query-gpu=memory.total} with graceful fallback to 0\,GB
\end{itemize}

\subsection{Classification Matrix}

Based on detected hardware, systems are classified into four tiers as shown in Table~\ref{tab:tiers}.

\begin{table}[htbp]
\caption{Hardware Classification and AI Model Selection Matrix}
\label{tab:tiers}
\centering
\renewcommand{\arraystretch}{1.2}
\begin{tabular}{@{}p{1.2cm}p{1.8cm}p{2.0cm}p{2.0cm}@{}}
\toprule
\textbf{Tier} & \textbf{Requirements} & \textbf{LLM Model} & \textbf{Optimizations} \\
\midrule
Low-End & RAM $<$ 8\,GB, No GPU & Qwen 1.5B & ZRAM $1.5\times$, Swappiness=10 \\
\addlinespace
Mid-Range & RAM 8--16\,GB, VRAM $<$ 8\,GB & Qwen 3B / Llama-3 8B & Thread cap, 8K context \\
\addlinespace
High-End & RAM 16--32\,GB, VRAM $\geq$ 8\,GB & Qwen 14B & GPU offload, unified memory \\
\addlinespace
Work-station & RAM $>$ 32\,GB, VRAM $\geq$ 12\,GB & Llama-3.3 70B & 100\% GPU, FlashAttention-2 \\
\bottomrule
\end{tabular}
\end{table}

\subsection{Performance Scoring}

A composite performance score quantifies system capability:
\begin{equation}
S = 3.0 \times V_{\text{GPU}} + 0.5 \times M_{\text{RAM}} + 0.2 \times N_{\text{cores}}
\end{equation}
where $V_{\text{GPU}}$ is GPU VRAM in GB, $M_{\text{RAM}}$ is system RAM in GB, and $N_{\text{cores}}$ is the CPU core count. GPU memory is weighted 6$\times$ higher than RAM per gigabyte, reflecting the disproportionate impact of VRAM on LLM inference throughput.

\subsection{GPU Offload Strategy}

For mid-range systems with partial GPU availability, the offload percentage is calculated as:
\begin{equation}
O_{\text{GPU}} = \min\left(100, \frac{V_{\text{GPU}}}{6.0} \times 100\right)\%
\end{equation}
This enables proportional model layer distribution between GPU and CPU memory, maximizing inference speed without exceeding VRAM capacity.

\subsection{First-Boot Auto-Tuning}

Upon initial boot, the \texttt{firstboot\_setup.py} service:
\begin{enumerate}
    \item Executes the hardware profiler
    \item Seeds recommended model configuration into SQLite
    \item Configures ZRAM with LZ4 compression at $1.5\times$ RAM capacity (low-end systems)
    \item Sets kernel swappiness to 5--10 to prevent LLM parameter page-outs
    \item Writes a sentinel file to prevent re-execution
\end{enumerate}

% ============================================================================
%                   VI. SECURITY ARCHITECTURE
% ============================================================================

\section{Defense-in-Depth Security Architecture}
\label{sec:security}

M.A.T.R.I.X. enforces security through four independent containment layers, ensuring that breach of any single layer does not compromise system integrity.

\subsection{Layer 1: AppArmor Mandatory Access Control}

A custom AppArmor profile (\texttt{matrix-sandbox-profile}) implements path-based MAC with the following policy:

\begin{itemize}
    \item \textbf{Executable whitelist}: Only \texttt{/usr/bin/\{bash, sh, ls, pwd, whoami, python3*\}} may execute
    \item \textbf{Filesystem confinement}: Read-write access restricted to \texttt{/var/lib/matrix/}, \texttt{/home/matrix/}, \texttt{/tmp/matrix\_build/}
    \item \textbf{Critical file protection}: Write access denied to \texttt{/etc/passwd}, \texttt{/etc/shadow}, \texttt{/etc/sudoers}, \texttt{/boot/}
    \item \textbf{Network restriction}: Raw and packet socket creation denied, preventing packet sniffing and IP spoofing
\end{itemize}

\subsection{Layer 2: Bubblewrap Namespace Isolation}

Interactive terminal tasks execute within Bubblewrap containers using comprehensive namespace isolation:

\begin{itemize}
    \item \texttt{--unshare-all}: Isolates PID, network, user, mount, IPC, and UTS namespaces
    \item System directories (\texttt{/usr}, \texttt{/lib}, \texttt{/bin}, \texttt{/sbin}) mounted read-only via \texttt{--ro-bind}
    \item Only \texttt{/var/lib/matrix} and \texttt{/home/matrix} mounted read-write
    \item Fresh \texttt{/tmp}, \texttt{/proc}, and \texttt{/dev} created within the container
\end{itemize}

\subsection{Layer 3: Systemd Service Hardening}

Each daemon service unit specifies eight security directives:

\begin{table}[htbp]
\caption{Systemd Security Hardening Directives}
\label{tab:systemd}
\centering
\renewcommand{\arraystretch}{1.15}
\begin{tabular}{@{}ll@{}}
\toprule
\textbf{Directive} & \textbf{Effect} \\
\midrule
\texttt{ProtectSystem=strict} & Read-only filesystem \\
\texttt{ProtectHome=yes} & Block home directory access \\
\texttt{PrivateTmp=yes} & Isolated /tmp namespace \\
\texttt{PrivateDevices=yes} & No device node access \\
\texttt{DevicePolicy=closed} & Block device creation \\
\texttt{CapabilityBoundingSet=} & Drop ALL capabilities \\
\texttt{NoNewPrivileges=true} & Prevent privilege escalation \\
\texttt{ReadWritePaths=...} & Explicit write whitelist \\
\bottomrule
\end{tabular}
\end{table}

\subsection{Layer 4: Application-Level Permission System}

The API daemon classifies all terminal commands into three tiers:

\begin{itemize}
    \item \textbf{Safe} (auto-execute): Commands matching a prefix whitelist (\texttt{ls}, \texttt{pwd}, \texttt{git status}, etc.), parsed with \texttt{shlex.split} and executed with \texttt{shell=False}
    \item \textbf{Confirm} (user approval required): Commands containing shell chaining characters (\texttt{;}, \texttt{\&}, \texttt{|}, \texttt{\`{}}, \texttt{\$})
    \item \textbf{Admin} (elevated operations): System-modifying commands routed to a pending approval queue with UUID action tracking
\end{itemize}

All executions are logged to \texttt{/var/log/matrix\_api\_audit.log} with timestamp, permission level, exit code, and output metadata.

% ============================================================================
%                    VII. USER INTERFACE
% ============================================================================

\section{Operator Control Interface}
\label{sec:ui}

The frontend is implemented as a single-page application (SPA) using React~19 with Vite~6 and the SWC compilation plugin. The interface simulates a complete Ubuntu desktop environment with the following components:

\subsection{Desktop Environment}

The desktop metaphor includes a top status bar (clock, WiFi, volume, battery), left dock sidebar with application launchers, desktop icon grid, and a window manager with close/minimize/maximize controls. Audio feedback is generated through the Web Audio API using oscillator synthesis.

\subsection{Application Modules}

Nine integrated application tabs provide comprehensive system control:

\begin{enumerate}
    \item \textbf{Terminal}: NLP-powered command parser supporting 20+ Linux commands with Ollama integration and offline fallback. Features a 4-step system pipeline visualization.
    
    \item \textbf{Build Factory}: 6-stage ISO build pipeline simulator with real-time log generation and agent status monitoring.
    
    \item \textbf{System Daemon}: Three-panel hardware telemetry dashboard displaying Ollama inference metrics, Wine/Proton syscall statistics, and host resource utilization with live/simulated mode detection.
    
    \item \textbf{Sandbox Manager}: Configuration panel for AppArmor profiles, cgroup limits, nftables rules, and escape corruption monitoring.
    
    \item \textbf{M.A.T.R.I.X. AI}: Voice-controlled AI assistant with a 3D holographic particle sphere using Fibonacci lattice distribution (120 particles with perspective projection). Integrates Web Speech API for recognition and synthesis, with 40+ voice command intent matchers.
    
    \item \textbf{Nautilus}: Simulated file explorer with breadcrumb navigation and file content preview.
    
    \item \textbf{Browser}: Web browser simulator with hardcoded sites and navigation history.
    
    \item \textbf{Game}: Canvas-based ``Firewall Bypass'' snake game with collision detection and progressive difficulty.
    
    \item \textbf{Drivers}: GPU (NVIDIA/Nouveau) and audio (PipeWire/ALSA) driver configuration with cross-tab communication via \texttt{CustomEvent}.
\end{enumerate}

\subsection{3D Holographic Visualization}

The AI assistant features a real-time 3D particle sphere rendered on HTML5 Canvas. Particles are distributed using the Fibonacci lattice algorithm:
\begin{equation}
\theta_i = \arccos(1 - 2i/N), \quad \phi_i = \pi(1 + \sqrt{5}) \cdot i
\end{equation}
where $N = 120$ particles. The sphere rotates continuously with perspective projection and changes color based on system state: cyan (dormant), orange (listening), violet (processing), and pulsing cyan (speaking). Microphone input drives radius pulsation through real-time FFT analysis via the Web Audio API \texttt{AnalyserNode}.

% ============================================================================
%                    VIII. ISO BUILD PIPELINE
% ============================================================================

\section{Automated ISO Build Pipeline}
\label{sec:build}

The ISO compilation follows a 13-stage pipeline:

\begin{enumerate}[label=\textbf{S\arabic*}:]
    \item \textbf{Clean}: Remove previous build artifacts
    \item \textbf{Validate}: Check for required tools (\texttt{debootstrap}, \texttt{mksquashfs}, \texttt{xorriso}, \texttt{grub-mkrescue})
    \item \textbf{Bootstrap}: \texttt{debootstrap --arch=amd64 bookworm} creates minimal Debian chroot
    \item \textbf{Copy}: Transfer all M.A.T.R.I.X. source files, service units, and profiles into chroot
    \item \textbf{Mount}: Bind virtual filesystems (\texttt{/proc}, \texttt{/sys}, \texttt{/dev})
    \item \textbf{Chroot}: Execute customization script installing 90+ packages, configuring Wine, Ollama, and the \texttt{matrix} user
    \item \textbf{Unmount}: Clean virtual filesystem binds
    \item \textbf{Extract}: Copy kernel (\texttt{vmlinuz}) and initrd from chroot
    \item \textbf{Cleanup}: Remove temporary build files
    \item \textbf{SquashFS}: Compress root filesystem with XZ algorithm
    \item \textbf{GRUB}: Generate bootloader configuration for BIOS/EFI dual boot
    \item \textbf{ISO}: Package using \texttt{xorriso} into bootable ISO image
    \item \textbf{Checksum}: Generate MD5 integrity hash
\end{enumerate}

Trap-based cleanup (\texttt{trap cleanup EXIT ERR INT TERM}) ensures that virtual filesystem mounts are properly unmounted even on build failures.

\subsection{CI/CD Multi-Agent Workflow}

Ten GitHub Actions workflows implement a cascading multi-agent pipeline:

\begin{equation*}
\text{Master} \rightarrow \text{Architect} \rightarrow \begin{cases} \text{Developer} \\ \text{UI/UX} \\ \text{Package} \end{cases} \rightarrow \text{Debug} \rightarrow \text{Build}
\end{equation*}

A Learning Agent activates on Debug/Build failures, and an Updater Agent runs on a daily cron schedule. The CI pipeline validates shell scripts (ShellCheck), React builds, and Python compilation on every push.

% ============================================================================
%                    IX. EXPERIMENTAL EVALUATION
% ============================================================================

\section{Experimental Evaluation}
\label{sec:results}

\subsection{System Metrics}

Table~\ref{tab:codebase} summarizes the complete codebase composition.

\begin{table}[htbp]
\caption{M.A.T.R.I.X. Codebase Composition}
\label{tab:codebase}
\centering
\renewcommand{\arraystretch}{1.15}
\begin{tabular}{@{}lrrr@{}}
\toprule
\textbf{Component} & \textbf{Files} & \textbf{Lines} & \textbf{Language} \\
\midrule
Core Daemons & 3 & 1,460 & Python \\
Scheduler & 2 & 988 & Python \\
ISO Build Scripts & 12 & 835 & Bash/Config \\
Frontend UI & 10 & 6,981 & JSX/CSS \\
Documentation & 9 & --- & Markdown \\
CI/CD Workflows & 10 & --- & YAML \\
\midrule
\textbf{Total Backend} & \textbf{17} & \textbf{3,283} & --- \\
\bottomrule
\end{tabular}
\end{table}

\subsection{Scheduling Performance}

The scheduler simulation suite (\texttt{scheduler\_prototype.py}) evaluates all four algorithms with identical workloads. Table~\ref{tab:scheduling} presents comparative results.

\begin{table}[htbp]
\caption{Scheduling Algorithm Comparison (5-Task Workload)}
\label{tab:scheduling}
\centering
\renewcommand{\arraystretch}{1.15}
\begin{tabular}{@{}lcccc@{}}
\toprule
\textbf{Metric} & \textbf{FIFO} & \textbf{SJF} & \textbf{RR} & \textbf{Priority} \\
\midrule
Avg. Wait (s) & 4.2 & 2.8 & 3.5 & 3.1 \\
Context Switches & 5 & 5 & 18 & 5 \\
Starvation Risk & No & High & No & Mitigated \\
Fairness & Low & Low & High & Medium \\
\bottomrule
\end{tabular}
\end{table}

SJF achieves the lowest average waiting time (2.8\,s), consistent with theoretical optimality \cite{silberschatz2018os}. Round Robin provides the highest fairness but incurs 3.6$\times$ more context switches. Priority with Aging achieves a balanced trade-off with near-SJF wait times while mitigating starvation through the urgency increment mechanism ($\Delta_{\text{age}} = 0.8$).

\subsection{LRU-K Cache Performance}

With cache capacity $=2$ and $K=2$:
\begin{itemize}
    \item \textbf{Cache hit latency}: 100\,ms (50\,ms save + 50\,ms restore)
    \item \textbf{Cache miss latency}: 200\,ms (50\,ms save + 150\,ms database load)
    \item \textbf{Latency reduction on hit}: 50\%
\end{itemize}

For workloads exhibiting temporal locality---where the same agents are frequently rescheduled---the LRU-2 policy maintains hit rates exceeding 60\%, reducing average context switch overhead by approximately 30\%.

\subsection{Boot Verification}

The headless QEMU validation script (\texttt{verify\_boot.sh}) confirms successful boot-to-service initialization:
\begin{itemize}
    \item Target services (\texttt{matrix-daemon}, \texttt{matrix-api}) register within 60\,s
    \item Login prompt appears confirming kernel and init system completion
    \item Serial console log monitoring at 2-second polling intervals
    \item Exit code 0/1 enables CI/CD gate integration
\end{itemize}

\subsection{Security Validation}

The sandbox test suite (\texttt{test\_sandbox.sh}) validates:
\begin{itemize}
    \item Bubblewrap read-only enforcement: write attempts to \texttt{/usr} return ``Read-only file system''
    \item AppArmor profile syntax: \texttt{apparmor\_parser -Q} validates profile correctness
    \item Path traversal prevention: \texttt{os.path.commonpath()} confinement blocks directory escape attempts with 403 responses
\end{itemize}

% ============================================================================
%                      X. FUTURE WORK
% ============================================================================

\section{Future Work and Upgrade Roadmap}
\label{sec:future}

The following enhancements are planned for future iterations:

\subsection{Near-Term Upgrades}

\begin{enumerate}
    \item \textbf{Local Vector RAG Memory}: Integration of ChromaDB or SQLite-vss for semantic search over user files, shell history, and notes. Local embeddings via \texttt{all-minilm-l6-v2} (45\,MB) will enable context-aware AI responses across sessions without cloud data transmission.
    
    \item \textbf{Neural Text-to-Speech}: Replacement of Web Speech API synthesis with Piper TTS (CPU-optimized, ONNX format) or Kokoro TTS (82M parameters) for natural, offline voice output via ALSA/PipeWire audio pipelines.
    
    \item \textbf{Whisper Speech-to-Text}: Integration of OpenAI Whisper \cite{radford2023whisper} (tiny/base models, 39--74M parameters) for offline speech recognition, replacing the browser-dependent Web Speech API.
    
    \item \textbf{CPU Pinning}: Reservation of one CPU core for the graphical interface with LLM threads pinned to remaining cores via \texttt{taskset}, preventing inference workloads from causing UI lag.
\end{enumerate}

\subsection{Medium-Term Enhancements}

\begin{enumerate}
    \item \textbf{Multi-Level Feedback Queue (MLFQ)}: Extension of the scheduler with MLFQ \cite{silberschatz2018os} to dynamically adjust task priority based on observed CPU burst patterns, combining the advantages of SJF and Round Robin.
    
    \item \textbf{Real-Time GPU Telemetry}: Integration of \texttt{nvidia-smi} and ROCm monitoring for live GPU utilization, VRAM consumption, and thermal tracking within the system daemon dashboard.
    
    \item \textbf{Container Orchestration}: Extension of Bubblewrap sandboxing to support multiple concurrent AI agent containers with resource quota enforcement via cgroups v2.
    
    \item \textbf{ARM64 Cross-Compilation}: Adaptation of the ISO build pipeline for ARM64 targets (Raspberry Pi 4/5, Orange Pi), enabling deployment on edge computing and IoT devices.
    
    \item \textbf{Wine/Proton Deep Integration}: Full DXVK/VKD3D pipeline integration for running Windows PE binaries with NTSync kernel-level NT synchronization for reduced syscall translation overhead.
\end{enumerate}

\subsection{Long-Term Vision}

\begin{enumerate}
    \item \textbf{Federated Learning}: Privacy-preserving model fine-tuning across M.A.T.R.I.X. instances using federated averaging \cite{mcmahan2017federated}, enabling collective improvement without centralizing user data.
    
    \item \textbf{Agentic OS Autonomy}: Evolution toward autonomous agent swarms capable of self-directed system administration, package management, and security patching through multi-agent reinforcement learning.
    
    \item \textbf{Custom Kernel Modules}: Development of kernel-space scheduling modules for sub-millisecond agent context switching, replacing the current userspace \texttt{asyncio} implementation.
    
    \item \textbf{Distributed Agent Mesh}: Peer-to-peer agent collaboration across multiple M.A.T.R.I.X. instances on local networks using encrypted gRPC channels for collaborative task execution.
\end{enumerate}

% ============================================================================
%                       XI. CONCLUSION
% ============================================================================

\section{Conclusion}
\label{sec:conclusion}

This paper presented M.A.T.R.I.X., a comprehensive proof-of-concept for an AI-first operating system that integrates local LLM inference, multi-policy agent scheduling, hardware-adaptive optimization, and defense-in-depth security within a custom Linux distribution. The system demonstrates that meaningful AI-OS integration is achievable using open-source tools and standard library components, without requiring proprietary hardware, cloud infrastructure, or external package dependencies.

The multi-policy scheduler with LRU-K attention caching provides flexible task management with 50\% latency reduction on cache hits. The four-layer security architecture ensures robust containment of AI agent operations. The hardware profiling engine enables automatic deployment across hardware configurations ranging from 4\,GB RAM laptops to multi-GPU workstations.

The 13-stage ISO build pipeline, combined with headless boot verification and a 10-workflow CI/CD agent chain, establishes a reproducible, testable deployment methodology. The React-based glassmorphic desktop interface demonstrates that premium user experiences can coexist with system-level AI integration.

M.A.T.R.I.X. establishes a foundation for privacy-preserving, self-configuring AI operating systems and invites the broader research community to explore the convergence of artificial intelligence, operating system design, and local-first computing paradigms.

% ============================================================================
%                        ACKNOWLEDGMENTS
% ============================================================================

\section*{Acknowledgments}

The author acknowledges the open-source communities behind Debian, Ollama, React, Llama.cpp, Qwen, AppArmor, and Bubblewrap, whose foundational contributions made this work possible.

% ============================================================================
%                         REFERENCES
% ============================================================================

\begin{thebibliography}{30}

\bibitem{openai2023gpt4}
OpenAI, ``GPT-4 Technical Report,'' \textit{arXiv preprint arXiv:2303.08774}, 2023.

\bibitem{google2024gemini}
Gemini Team, ``Gemini: A Family of Highly Capable Multimodal Models,'' \textit{arXiv preprint arXiv:2312.11805}, 2024.

\bibitem{microsoft2024copilot}
Microsoft, ``Microsoft Copilot: Your AI companion,'' 2024. [Online]. Available: \url{https://copilot.microsoft.com}

\bibitem{meta2024llama3}
Meta AI, ``Llama 3: Open Foundation and Fine-Tuned Chat Models,'' \textit{Meta AI Blog}, 2024.

\bibitem{qwen2024qwen25}
Qwen Team, ``Qwen2.5: A Party of Foundation Models,'' \textit{arXiv preprint arXiv:2412.15115}, 2024.

\bibitem{microsoft2024phi3}
Microsoft Research, ``Phi-3 Technical Report: A Highly Capable Language Model Locally on Your Phone,'' \textit{arXiv preprint arXiv:2404.14219}, 2024.

\bibitem{ollama2024}
Ollama, ``Ollama: Get up and running with large language models,'' 2024. [Online]. Available: \url{https://ollama.ai}

\bibitem{llamacpp2024}
G. Gerganov, ``llama.cpp: Inference of Meta's LLaMA model in pure C/C++,'' GitHub, 2024.

\bibitem{shortliffe1976mycin}
E.~H. Shortliffe, \textit{Computer-Based Medical Consultations: MYCIN}. Elsevier, 1976.

\bibitem{giarratano2005expert}
J.~C. Giarratano and G.~D. Riley, \textit{Expert Systems: Principles and Programming}, 4th ed. Course Technology, 2005.

\bibitem{microsoft2020cortana}
Microsoft, ``Cortana, your personal productivity assistant,'' 2020.

\bibitem{apple2024siri}
Apple Inc., ``Siri --- Apple,'' 2024. [Online]. Available: \url{https://www.apple.com/siri/}

\bibitem{apple2024intelligence}
Apple Inc., ``Apple Intelligence,'' 2024. [Online]. Available: \url{https://www.apple.com/apple-intelligence/}

\bibitem{gpt4all2024}
Nomic AI, ``GPT4All: Run Large Language Models Locally,'' 2024.

\bibitem{lmstudio2024}
LM Studio, ``Discover, download, and run local LLMs,'' 2024. [Online]. Available: \url{https://lmstudio.ai}

\bibitem{jan2024}
Jan.ai, ``Jan: Open-source ChatGPT alternative,'' 2024. [Online]. Available: \url{https://jan.ai}

\bibitem{nixos2024}
NixOS Foundation, ``NixOS: The Purely Functional Linux Distribution,'' 2024.

\bibitem{silberschatz2018os}
A.~Silberschatz, P.~B. Galvin, and G.~Gagne, \textit{Operating System Concepts}, 10th ed. Wiley, 2018.

\bibitem{tanenbaum2015modern}
A.~S. Tanenbaum and H.~Bos, \textit{Modern Operating Systems}, 4th ed. Pearson, 2015.

\bibitem{stallings2018os}
W.~Stallings, \textit{Operating Systems: Internals and Design Principles}, 9th ed. Pearson, 2018.

\bibitem{oneil1993lruk}
E.~J. O'Neil, P.~E. O'Neil, and G.~Weikum, ``The LRU-K Page Replacement Algorithm for Database Disk Buffering,'' in \textit{Proc. ACM SIGMOD}, 1993, pp.~297--306.

\bibitem{apparmor2024}
AppArmor Project, ``AppArmor Security Profiles,'' 2024. [Online]. Available: \url{https://apparmor.net}

\bibitem{bubblewrap2024}
Project Atomic, ``Bubblewrap: Unprivileged sandboxing tool,'' GitHub, 2024.

\bibitem{systemd2024}
systemd Project, ``systemd System and Service Manager,'' 2024. [Online]. Available: \url{https://systemd.io}

\bibitem{flatpak2024}
Flatpak Team, ``Flatpak---the future of application distribution,'' 2024.

\bibitem{snap2024}
Canonical, ``Snapcraft: The app store for Linux,'' 2024.

\bibitem{firejail2024}
Firejail Project, ``Firejail Security Sandbox,'' GitHub, 2024.

\bibitem{debianlive2024}
Debian Project, ``Debian Live Manual,'' 2024.

\bibitem{lfs2024}
G.~Beekmans, ``Linux From Scratch,'' Version 12.1, 2024.

\bibitem{debootstrap2024}
Debian Project, ``debootstrap: Bootstrap a basic Debian system,'' 2024.

\bibitem{squashfs2024}
P.~Lougher, ``SquashFS: A compressed read-only filesystem for Linux,'' 2024.

\bibitem{debianpreseed2024}
Debian Project, ``Automating the installation using preseeding,'' \textit{Debian Installation Guide}, 2024.

\bibitem{qemu2024}
QEMU Project, ``QEMU: A generic and open source machine emulator and virtualizer,'' 2024.

\bibitem{radford2023whisper}
A.~Radford \textit{et al.}, ``Robust Speech Recognition via Large-Scale Weak Supervision,'' in \textit{Proc. ICML}, 2023.

\bibitem{mcmahan2017federated}
B.~McMahan \textit{et al.}, ``Communication-Efficient Learning of Deep Networks from Decentralized Data,'' in \textit{Proc. AISTATS}, 2017, pp.~1273--1282.

\end{thebibliography}

\balance

\end{document}
